% =========================================================================
\section{Impedancia, transformada de Laplace y respuesta en frecuencia}\label{sec:laplace}
% =========================================================================
El circuito~3 del guion contiene un condensador, cuya relación tensión--corriente es
diferencial. Esta sección reúne las tres ideas que permiten tratarlo con la misma
álgebra que una resistencia: la impedancia compleja, la transformada de Laplace y el
diagrama de Bode~\cite{TEORIA-LAPLACE,TEORIA-ADICIONAL}.

% -------------------------------------------------------------------------
\subsection{Régimen sinusoidal e impedancia compleja}\label{sec:impedancia}
% -------------------------------------------------------------------------
En un circuito lineal excitado por una sinusoide de pulsación $\omega=2\pi f$, todas las
tensiones y corrientes son, en régimen permanente, sinusoides de la misma pulsación. Cada una
queda determinada por su amplitud y su fase, que se agrupan en un número complejo, el \emph{fasor}:
\begin{equation}\label{eq:fasor}
  v(t)=V_p\cos(\omega t+\varphi)=\Real\bigl\{\mathbf{V}\,\eul^{\jj\omega t}\bigr\},\qquad
  \mathbf{V}=V_p\,\eul^{\jj\varphi}\in\C .
\end{equation}
Como $\frac{\dif}{\dif t}\eul^{\jj\omega t}=\jj\omega\,\eul^{\jj\omega t}$, \emph{derivar en el tiempo equivale a
multiplicar el fasor por $\jj\omega$}.

\begin{definicion}[Impedancia]\label{def:impedancia}
La impedancia de un dipolo lineal es el cociente de los fasores de su tensión y su corriente,
$Z(\jj\omega)=\mathbf{V}/\mathbf{I}$. Su módulo es el cociente de amplitudes y su argumento, el
desfase de la tensión respecto de la corriente.
\end{definicion}

\begin{proposicion}[Impedancias elementales]\label{prop:impedancias}
\begin{equation}\label{eq:impedancias}
  Z_R=R,\qquad Z_C=\frac{1}{\jj\omega C},\qquad Z_L=\jj\omega L .
\end{equation}
\end{proposicion}
\begin{proof}
Ley de Ohm, $v=Ri\Rightarrow\mathbf{V}=R\mathbf{I}$. En el condensador, $i=C\,\dif v/\dif t\Rightarrow
\mathbf{I}=\jj\omega C\,\mathbf{V}$. En la bobina, $v=L\,\dif i/\dif t\Rightarrow\mathbf{V}=\jj\omega L\,\mathbf{I}$.
\end{proof}

Las leyes de Kirchhoff son lineales, de modo que se cumplen también para los fasores; en
consecuencia, las impedancias se asocian como resistencias: en serie se suman,
$Z=Z_1+Z_2$, y en paralelo se suman sus inversas, $Z=Z_1Z_2/(Z_1+Z_2)$. En particular,
\begin{equation}\label{eq:RC-paralelo}
  Z_{R\parallel C}=\frac{R\cdot\frac{1}{\jj\omega C}}{R+\frac{1}{\jj\omega C}}=\frac{R}{1+\jj\omega RC} .
\end{equation}
El condensador es un circuito abierto en continua ($\abs{Z_C}\to\infty$ si $\omega\to0$) y un
cortocircuito a alta frecuencia ($\abs{Z_C}\to0$ si $\omega\to\infty$).

% -------------------------------------------------------------------------
\subsection{Transformada de Laplace y función de transferencia}\label{sec:transformada}
% -------------------------------------------------------------------------
Los fasores solo describen el régimen sinusoidal permanente. La transformada de Laplace
generaliza la idea a señales arbitrarias (cuadradas, triangulares, escalones), que son las que
se usan en el laboratorio para comprobar la integración y la derivación.

\begin{definicion}[Transformada de Laplace]\label{def:laplace}
Para una señal $f(t)$ definida en $t\ge0$, su transformada de Laplace es
\begin{equation}\label{eq:laplace}
  F(s)=\Lt\{f\}(s)=\int_{0}^{\infty}f(t)\,\eul^{-st}\,\dif t,\qquad s=\sigma+\jj\omega\in\C,
\end{equation}
definida en el semiplano $\Real s>\sigma_0$ en que la integral converge.
\end{definicion}

\begin{proposicion}[Propiedades que usaremos]\label{prop:laplace}
Si $f\leftrightarrow F$ y $g\leftrightarrow G$, con $a,b$ constantes:
\begin{gather}
  af(t)+bg(t)\;\longleftrightarrow\;aF(s)+bG(s),\label{eq:L-lineal}\\
  \frac{\dif f}{\dif t}\;\longleftrightarrow\;sF(s)-f(0),\qquad
  \int_{0}^{t}f(\tau)\,\dif\tau\;\longleftrightarrow\;\frac{F(s)}{s}.\label{eq:L-derivada}
\end{gather}
\end{proposicion}
\begin{proof}
La linealidad es la de la integral. Para la derivada se integra por partes:
$\int_0^\infty f'\eul^{-st}\dif t=\bigl[f\eul^{-st}\bigr]_0^\infty+s\int_0^\infty f\eul^{-st}\dif t=sF(s)-f(0)$.
Si $g(t)=\int_0^t f$, entonces $g'=f$ y $g(0)=0$, de modo que $F=sG$.
\end{proof}

Con condiciones iniciales nulas, \cref{eq:L-derivada} dice que \emph{derivar equivale a
multiplicar por $s$ e integrar, a dividir por $s$}. Aplicado al condensador,
$i=C\,\dif v/\dif t\Rightarrow I(s)=sC\,V(s)$, y se recupera \cref{eq:impedancias} con $\jj\omega\to s$:
\begin{equation}\label{eq:Zs}
  Z_R(s)=R,\qquad Z_C(s)=\frac{1}{sC},\qquad Z_L(s)=sL .
\end{equation}
Así, un circuito con condensadores se resuelve, en el dominio $s$, exactamente como uno resistivo.

\begin{definicion}[Función de transferencia]\label{def:transferencia}
Para un circuito lineal con condiciones iniciales nulas, la función de transferencia es el cociente
\begin{equation}\label{eq:transferencia}
  H(s)=\frac{V_{\mathrm{out}}(s)}{V_{\mathrm{in}}(s)} .
\end{equation}
Su restricción al eje imaginario, $H(\jj\omega)$, es la \emph{respuesta en frecuencia}.
\end{definicion}

\begin{proposicion}[Respuesta a una sinusoide]\label{prop:sinusoide}
Si $H(s)$ es estable (polos en $\Real s<0$) y $\vin(t)=V_p\cos\omega t$, la salida en régimen
permanente es
\begin{equation}\label{eq:respuesta-sinusoide}
  \vout(t)=\abs{H(\jj\omega)}\,V_p\cos\bigl(\omega t+\arg H(\jj\omega)\bigr).
\end{equation}
\end{proposicion}
\begin{proof}
La exponencial $\eul^{\jj\omega t}$ es autofunción de todo sistema lineal e invariante: si
$\vin=\eul^{\jj\omega t}$, la salida forzada es $H(\jj\omega)\eul^{\jj\omega t}$ (los términos asociados a los
polos de $H$ decaen como $\eul^{\Real(p)t}\to0$). Tomando la parte real, por linealidad, se obtiene el
resultado.
\end{proof}
La ganancia medida con el osciloscopio para una sinusoide de frecuencia $f$ es, por tanto,
$\abs{A_V}=\abs{H(\jj 2\pi f)}$, y el desfase entre los canales es $\arg H(\jj2\pi f)$.

% -------------------------------------------------------------------------
\subsection{Diagrama de Bode y red de constante temporal única}\label{sec:bode}
% -------------------------------------------------------------------------
\begin{definicion}[Diagrama de Bode]\label{def:bode}
Es la representación de $\abs{H(\jj\omega)}$ en decibelios, $20\log_{10}\abs{H}$, y de la fase
$\arg H(\jj\omega)$, ambas frente a la frecuencia en escala logarítmica. Una \emph{década} es un
factor $10$ en frecuencia y una \emph{octava}, un factor $2$.
\end{definicion}

\begin{proposicion}[Red de constante temporal única paso bajo]\label{prop:STC}
Sea
\begin{equation}\label{eq:STC}
  H(s)=\frac{K}{1+s/\omega_0},\qquad \omega_0=\frac1\tau .
\end{equation}
Entonces
\begin{equation}\label{eq:STC-modulo}
  \abs{H(\jj\omega)}=\frac{\abs{K}}{\sqrt{1+(\omega/\omega_0)^{2}}},\qquad
  \arg H(\jj\omega)=\arg K-\arctg\frac{\omega}{\omega_0},
\end{equation}
y el diagrama de Bode tiene dos asíntotas que se cortan en $\omega_0$ (\cref{fig:bode-stc}):
\begin{equation}\label{eq:STC-asintotas}
  \omega\ll\omega_0:\ \abs{H}\simeq\abs{K};\qquad
  \omega\gg\omega_0:\ \abs{H}\simeq\abs{K}\,\frac{\omega_0}{\omega},
  \ \text{pendiente }-20\dB/\text{década}\ (-6\dB/\text{octava}).
\end{equation}
En $\omega=\omega_0$ el módulo es $\abs{K}/\sqrt2$, es decir, $3\dB$ por debajo de la meseta, y la fase
se ha retrasado $45^{\circ}$: $\omega_0$ es la \emph{frecuencia de corte} a $-3\dB$.
\end{proposicion}
\begin{proof}
Con $s=\jj\omega$, $\abs{1+\jj x}=\sqrt{1+x^{2}}$ y $\arg(1+\jj x)=\arctg x$, con $x=\omega/\omega_0$. Si $x\ll1$,
$\sqrt{1+x^{2}}\simeq1$; si $x\gg1$, $\sqrt{1+x^{2}}\simeq x$, y $20\log_{10}(\abs{K}/x)$ disminuye
$20\dB$ cada vez que $x$ se multiplica por $10$ ($20\log_{10}2\simeq6\dB$ por octava). En $x=1$,
$20\log_{10}(1/\sqrt2)=-3{,}01\dB$ y $\arctg1=45^{\circ}$.
\end{proof}

Análogamente, $H(s)=K\,(s/\omega_0)/(1+s/\omega_0)$ es un \emph{paso alto}: crece a
$+20\dB$/década por debajo de $\omega_0$ y es plano por encima. Los dos casos límite
$H=K\omega_0/s$ (integrador, recta de $-20\dB$/década en todo el eje) y $H=Ks/\omega_0$
(derivador, $+20\dB$/década) son las asíntotas de alta y de baja frecuencia, respectivamente,
de estas dos redes.

% Figura: diagrama de Bode normalizado de una red STC paso bajo (K=1).
\begin{figure}[ht]
\centering
\begin{tikzpicture}
  \begin{semilogxaxis}[tema, name=mod, width=0.47\textwidth, height=5.2cm,
      xmin=0.01, xmax=100, ymin=-42, ymax=6, ytick={-40,-30,-20,-10,-3,0},
      xlabel={$\omega/\omega_0$}, ylabel={$\abs{H}/\abs{K}$ (dB)},
      log basis x=10, title style={font=\small}, title={(a) Módulo}]
    \addplot[asintota, domain=0.01:1] {0};
    \addplot[asintota, domain=1:100, samples=2] {-20*log10(x)};
    \addplot[azulacero, line width=1.1pt, domain=0.01:100, samples=200] {-10*log10(1+x^2)};
    \addplot[only marks, mark=*, mark size=1.6pt, rojoteo] coordinates {(1,-3.01)};
    \node[font=\scriptsize, anchor=south west, text=rojoteo] at (axis cs:1.15,-3) {$-3\dB$};
    \node[font=\scriptsize, anchor=south west, text=grispizarra, rotate=-33] at (axis cs:4,-11) {$-20\dB$/década};
  \end{semilogxaxis}
  \begin{semilogxaxis}[tema, at={(mod.east)}, anchor=west, xshift=1.9cm,
      width=0.47\textwidth, height=5.2cm,
      xmin=0.01, xmax=100, ymin=-95, ymax=5, ytick={-90,-45,0},
      xlabel={$\omega/\omega_0$}, ylabel={$\arg H-\arg K$ (grados)},
      log basis x=10, title style={font=\small}, title={(b) Fase}]
    \addplot[asintota, domain=0.01:0.1] {0};
    \addplot[asintota, domain=0.1:10, samples=2] {-45-45*log10(x)};
    \addplot[asintota, domain=10:100] {-90};
    \addplot[azulacero, line width=1.1pt, domain=0.01:100, samples=200] {-atan(x)};
    \addplot[only marks, mark=*, mark size=1.6pt, rojoteo] coordinates {(1,-45)};
    \node[font=\scriptsize, anchor=south west, text=rojoteo] at (axis cs:1.15,-44) {$-45^{\circ}$};
  \end{semilogxaxis}
\end{tikzpicture}
\caption{Diagrama de Bode de la red de constante temporal única paso bajo
$H(s)=K/(1+s/\omega_0)$, normalizado. En trazo discontinuo, las asíntotas de
\cref{eq:STC-asintotas} (en la fase, la aproximación usual de $-45^{\circ}$/década entre
$\omega_0/10$ y $10\,\omega_0$). En la frecuencia de corte el módulo cae $3\dB$ y la fase se retrasa
$45^{\circ}$.}
\label{fig:bode-stc}
\end{figure}

